\begin{table*}[tb!]
\centering
\resizebox{1.0\textwidth}{!}{% <-
\begin{tabular}{l|ccc|ccc}
\hline
\multirow{2}{*}{\textbf{Model}} & \multicolumn{3}{c|}{\textbf{SNIPS-NLU}} & \multicolumn{3}{c}{\textbf{ATIS}} \\ \cline{2-7} 
                      & Slot (F1) &Intent (Acc) & Overall (Acc) 
                      & Slot (F1) &Intent (Acc) & Overall (Acc)           \\ \hline
CNN TriCRF \cite{xu2013convolutional} 	&-		&-		&-	&0.944   &-	&-	\\
Joint Seq. \cite{hakkani2016multi}  	&0.873	&0.969	&0.732	&0.942   &0.926	&0.807	\\
Attention BiRNN \cite{liu2016attention} &0.878	&0.967	&0.741	&0.942   &0.911	&0.789	\\
Slot-Gated Full Atten. \cite{goo2018slot}	&0.888	&0.970	&0.755	&0.948   &0.936	&0.822	\\
DR-AGG \cite{gong2018information}       & -		&0.966  	& -		& -   	 &0.914 & -	\\
IntentCapsNet \cite{xia2018zero}  		& -		&0.974	& -		& -   	 &0.948 & -	\\
\hline
{\ModelName} &\textbf{0.918}&0.973 &\textbf{0.809}   &\textbf{0.952} &\textbf{0.950} & \textbf{0.834}\\
{\ModelName} w/o Intent Detection &0.902	& -	& -	&{0.948} &-	& -	\\ 
{\ModelName} w/o Joint Training     &0.902	& \textbf{0.977}	&0.804	&0.948  &0.847	&0.743\\ 
\hline
\end{tabular}
}
\vspace{0.15in}
\caption{Slot filling and intention detection results using {\ModelName} on two datasets.}
\label{tab::overall}
\end{table*}
\begin{figure*}[tb!]
    \centering
    \includegraphics[width=\linewidth]{figures/benchmark.pdf}
    \vspace{0.05in}
    \caption{Stratified 5-fold cross validation for benchmarking with existing NLU services on SNIPS-NLU dataset. Black bars indicate the standard deviation.}\label{fig::benchmark}
\end{figure*}

We use the validation data to choose hyperparameters.
For both datasets, we randomly initialize word embeddings using Xavier initializer and let them train with the model.
In the loss function, the down-weighting coefficient $\lambda$ is 0.5, margins $m^ +$ and $m^ -$ are set to 0.8 and 0.2 for all the existing intents. $\alpha$ is set as 0.1.
RMSProp optimizer \citep{tieleman2012lecture} is used to minimize the loss.  To alleviate over-fitting, we add the dropout to the LSTM layer with a  dropout rate of 0.2.

\section{Results}
\noindent\textbf{Quantitative Evaluation} The intent detection results on two datasets are reported in Table \ref{tab::overall}, where the proposed capsule-based model performs consistently better than current learning schemes for joint slot filling and intent detection, as well as capsule-based neural network models that only focuses on intent detection. These results demonstrate the novelty of the proposed capsule-based model {\ModelName} in jointly modeling the hierarchical relationships among words, slots and intents via the dynamic routing between capsules.

Also, we benchmark the intent detection performance of the proposed model with existing natural language understanding services\footnote{https://www.slideshare.net/KonstantinSavenkov/nlu-intent-detection-benchmark-by-intento-august-2017} in Figure \ref{fig::benchmark}.
Since the original data split is not available, we report the results with stratified 5-fold cross validation. From Figure \ref{fig::benchmark} we can see that the proposed model {\ModelName} is highly competitive with off-the-shelf systems that are available to use. Note that, our model archieves the performance without using pre-trained word representations: the word embeddings are simply trained from scratch.

\noindent\textbf{Ablation Study}
To investigate the effectiveness of {\ModelName} in joint slot filling and intent detection, we also report ablation test results in Table \ref{tab::overall}. ``w/o Intent Detection'' is the model without intent detection: only a dynamic routing is performed between {\FirstCapsule} and {\SecondCapsule} for the slot filling task, where we minimize $\mathcal{L}_{slot}$ during training; ``w/o Joint Training'' adopts a two-stage training where the model is first trained for slot filling by minimizing $\mathcal{L}_{slot}$, and then use the fixed slot representations to train for the intent detection task which minimizes $\mathcal{L}_{intent}$. From the lower part of Table \ref{tab::overall} we can see that by using a capsule-based hierarchical modeling between words and slots, the model {\ModelName} w/o Intent Detection is already able to outperform current alternatives on slot filling that adopt a sequential labeling schema. The joint training of slot filling and intent detection is able to give each subtask further improvements when the model parameters are updated jointly.

\noindent\textbf{Visualizing Agreement Values between Capsule Layers}
Thanks to the dynamic routing-by-agreement schema, the dynamically learned agreement values between different capsule layers naturally reflect how low-level features are collectively aggregated into high-level ones for each input utterance. In this section, we harness the intepretability of the proposed capsule-based model via hierarchical modeling and provide case studies and visualizations.

\textbf{Between WordCaps and SlotCaps}
First we study the agreement value $c_{kt}$ between the $t$-th word in the WordCaps and the $k$-th slot capsule in SlotCaps.
As shown in Figure \ref{fig:w2s_dist}, we observe that the dynamic routing-by-agreement is able to converge to an agreement quickly after the first iteration (shown in blue bars). It is able to assign a confident probability assignment close to 0 or 1. After the second iteration (shown in orange bars), the model is more certain about the routing decisions: probabilities are more leaning towards 0 or 1 as the model is confident about routing a word in WordCaps to its most appropriate slot in SlotCaps.

\begin{figure}[h]
    \centering
    \includegraphics[width=\linewidth]{figures/w2s_dist}
    \caption{The distribution of all agreement values between WordCaps and SlotCaps on the test split of SNIPS-NLU dataset. Blue: the distribution of values after the first iteration. Yellow: the distribution after the second iteration.}
    \label{fig:w2s_dist}
\end{figure}

However, we do find that when unseen slot values like new object names emerge in utterances like \texttt{show me the movie operetta for the theatre organ} with an intent of \texttt{SearchCreativeWork}, the iterative dynamic routing process would be even more appealing.
Figure \ref{fig:w2s} shows the agreement values learned by dynamic routing-by-agreement. Since the dynamic routing-by-agreement is an iterative process controlled by the variable $iter_{slot}$, we show the agreement values after the first iteration in the left part of Figure \ref{fig:w2s}, and the values after the second iteration in the right part.
 
\begin{figure}[t!]
    \centering
    \includegraphics[width=\linewidth]{figures/w2s.pdf}
    \caption{The learned agreement values between WordCaps (x-axis) and SlotCaps (y-axis). A sample from the test split of SNIPS-NLU dataset is shown (Left: after the fist routing iteration. Right: after the second iteration). Due to space limitations, only part of slots (7/72) are shown on the y-axis.}
    \label{fig:w2s}
\end{figure}
\begin{figure}[t!]
    \centering
    \includegraphics[width=\linewidth]{figures/s2i.pdf}
    \caption{The learned agreement values between SlotCaps (y-axis) and IntentCaps (x-axis). Left: after the first iteration. Right: after the second iteration. The same sample utterance used in Figure 5 is used here.}
    \label{fig:s2i}
    \vspace{-0.3in}
\end{figure}

From the left part of Figure \ref{fig:w2s}, we can see that after the first iteration, the model considers the word \texttt{operetta} itself alone is likely to be an object name, probably because the following word \texttt{for} is usually a context word being annotated as \texttt{O}. Thus it tends to route word \texttt{for} to both the slot \texttt{O} and the slot \texttt{I-object\_name}. However, from the right part of Figure \ref{fig:w2s} we can see that after the second iteration, the dynamic routing found an agreement and is more certain to have \texttt{operetta for the theatre organ} as a whole for the slot \texttt{B-object\_name} and \texttt{I-object\_name}.

\textbf{Between SlotCaps and IntentCaps} Similarly, we visualize the agreement values between each slot capsule in SlotCaps and each intent capsule in IntentCaps.
The left part of Figure \ref{fig:s2i} shows that after the first iteration, since the model is not able to correctly recognize \texttt{operetta for the theatre organ} as a whole, only the context slot \texttt{O} (correspond to the word \texttt{show me the}) and \texttt{B-object\_name} (correspond to the word \texttt{operetta}) contribute significantly to the final intent capsule. From the right part of Figure \ref{fig:s2i}, we found that with the word  \texttt{operetta for the theatre organ} being recognized in the lower capsule, the slots \texttt{I-object\_name} and \texttt{B-object\_type} contribute more to the correct intent capsule \texttt{SearchCreativeWork}, when comparing with other routing alternatives to other intent capsules.